\documentclass[11pt]{article}

\usepackage[margin=1in]{geometry}
\usepackage{amsmath,amssymb,amsthm}
\usepackage{graphicx}
\usepackage[hidelinks]{hyperref}

\newtheorem{theorem}{Theorem}
\newtheorem{lemma}{Lemma}

\newcommand{\eps}{\varepsilon}
\newcommand{\teneps}{\widehat{\otimes}_{\varepsilon}}
\newcommand{\stens}{\widehat{\otimes}_{\varepsilon,s,N}}
\newcommand{\Ball}{B}
\newcommand{\Sphere}{S}
\newcommand{\R}{\mathbb R}
\newcommand{\N}{\mathbb N}

\begin{document}

\begin{center}
{\Large ASQ Passes to Symmetric Injective Tensor Powers}\\[4pt]
{\large Candidate full solution to Question 4.3 of arXiv:1602.07090}
\end{center}

\section*{Source and Open Question}

The source paper is Johann Langemets, Vegard Lima, and Abraham Rueda Zoca,
\emph{Almost square and octahedral norms in tensor products of Banach spaces},
arXiv:1602.07090 \cite{LLRZ}. In Section 4, Question 4.3, the authors ask:

\begin{quote}
Let \(X\) be a Banach space. Is
\(\widehat{\otimes}_{\varepsilon,s,N}X\) ASQ for each \(N\in\mathbb N\)
whenever \(X\) is ASQ?
\end{quote}

The question appears on page 14 of the source PDF:

\begin{center}
\includegraphics[width=0.94\linewidth]{figures/open_problem_crop.png}
\end{center}

\section*{Definitions}

For a Banach space \(X\), the \(N\)-fold injective symmetric tensor product
\(\widehat{\otimes}_{\varepsilon,s,N}X\) is the completion of the linear span
of tensors \(x^N=x\otimes\cdots\otimes x\) under the norm
\[
 \left\|\sum_j \lambda_j x_j^N\right\|
 =
 \sup_{x^*\in \Ball_{X^*}}
 \left|\sum_j \lambda_j x^*(x_j)^N\right|.
\]

We use the standard equivalent finite-set form of ASQ: a Banach space \(Z\)
is ASQ if and only if, for every \(z_1,\ldots,z_m\in \Sphere_Z\) and every
\(\eps>0\), there is \(w\in \Sphere_Z\) such that
\[
  \|z_i\pm w\|\le 1+\eps \qquad (1\le i\le m).
\]
We also use the finite-dimensional ASQ extension theorem of
Abrahamsen, Langemets, and Lima \cite[Theorem 2.4]{ALL}: if \(X\) is ASQ,
then for every finite-dimensional subspace \(E\subset X\) and every
\(\delta>0\), there is \(y\in \Sphere_X\) such that
\[
  \|e+\lambda y\|\le (1+\delta)\max\{\|e\|,|\lambda|\}
  \qquad(e\in E,\ \lambda\in\R).
\]

\section*{Proof Intuition}

The non-symmetric injective tensor proof in the source paper works because
an ASQ space contains, near any finite-dimensional piece, a new direction that
behaves almost like an \(\ell_\infty\)-summand. The same idea works for the
symmetric injective norm once one dualizes the estimate.

Indeed, if \(E+\R y\) is almost \(E\oplus_\infty \R\), then every
\(x^*\in \Ball_{X^*}\) satisfies
\[
  \|x^*|_E\|+|x^*(y)|\le 1+\delta.
\]
For a finite symmetric tensor supported in \(E\), its value at \(x^*\) is
bounded by a constant times \(\|x^*|_E\|^N\), while the perturbing tensor
\(y^N\) contributes \(|x^*(y)|^N\). The scalar inequality
\[
  r^N+s^N\le (r+s)^N
\]
then gives exactly the required ASQ estimate.

\begin{lemma}[dual estimate]\label{lem:dual}
Let \(E\subset X\) be finite-dimensional and let \(y\in \Sphere_X\) satisfy
\[
  \|e+\lambda y\|\le (1+\delta)\max\{\|e\|,|\lambda|\}
  \qquad(e\in E,\ \lambda\in\R).
\]
Then every \(x^*\in \Ball_{X^*}\) satisfies
\[
  \|x^*|_E\|+|x^*(y)|\le 1+\delta .
\]
\end{lemma}

\begin{proof}
Define \(T:E\oplus_\infty \R\to X\) by \(T(e,\lambda)=e+\lambda y\).
The hypothesis says that \(\|T\|\le 1+\delta\). Hence
\(x^*\circ T\) has norm at most \(1+\delta\) on \(E\oplus_\infty \R\).
Since \((E\oplus_\infty\R)^*=E^*\oplus_1\R\), this norm is precisely
\(\|x^*|_E\|+|x^*(y)|\).
\end{proof}

\begin{theorem}
Let \(X\) be an ASQ Banach space and let \(N\in\N\). Then
\(\widehat{\otimes}_{\varepsilon,s,N}X\) is ASQ.
\end{theorem}

\begin{proof}
Put \(Z=\widehat{\otimes}_{\varepsilon,s,N}X\). Fix
\(u_1,\ldots,u_m\in \Sphere_Z\) and \(\eps>0\). Choose
\(\eta>0\) and \(\delta>0\) so small that
\[
  \eta+(1+\eta)(1+\delta)^N\le 1+\eps .
\]
For each \(i\), choose a finite symmetric tensor
\[
  p_i=\sum_{j=1}^{k_i}\lambda_{ij}x_{ij}^N
\]
with \(\|u_i-p_i\|<\eta\). Let
\[
  E=\operatorname{span}\{x_{ij}:1\le i\le m,\ 1\le j\le k_i\}.
\]
By the ASQ finite-dimensional extension theorem, choose \(y\in \Sphere_X\)
such that
\[
  \|e+\lambda y\|\le (1+\delta)\max\{\|e\|,|\lambda|\}
  \qquad(e\in E,\ \lambda\in\R).
\]
Set \(v=y^N\in Z\). Since \(\|y\|=1\), Hahn--Banach gives
\(\|v\|=1\).

Fix \(i\) and \(x^*\in \Ball_{X^*}\). Write
\[
  r=\|x^*|_E\|,\qquad s=|x^*(y)|.
\]
By Lemma \ref{lem:dual}, \(r+s\le 1+\delta\). The tensor \(p_i\) is
supported in \(E\), so by Hahn--Banach applied to \(E\),
\[
 \left|\sum_{j=1}^{k_i}\lambda_{ij}x^*(x_{ij})^N\right|
 \le \|p_i\|\, r^N.
\]
Also \(\|p_i\|\le \|u_i\|+\eta=1+\eta\). Therefore, for either choice of
sign,
\[
\begin{aligned}
 \left|p_i(x^*)\pm x^*(y)^N\right|
 &\le \|p_i\|r^N+s^N\\
 &\le (1+\eta)(r^N+s^N)\\
 &\le (1+\eta)(r+s)^N\\
 &\le (1+\eta)(1+\delta)^N .
\end{aligned}
\]
Taking the supremum over \(x^*\in\Ball_{X^*}\) yields
\[
  \|p_i\pm v\|\le (1+\eta)(1+\delta)^N.
\]
Consequently,
\[
  \|u_i\pm v\|
  \le \|u_i-p_i\|+\|p_i\pm v\|
  < \eta+(1+\eta)(1+\delta)^N
  \le 1+\eps .
\]
This holds for every \(i=1,\ldots,m\), so \(Z\) is ASQ.
\end{proof}

\section*{Verification Notes}

The proof uses no computation. The included script
\texttt{code/scalar\_bound\_verifier.py} only audits the scalar inequality
\(r^N+s^N\le(r+s)^N\) on grids and random samples; the analytic reason is
the binomial theorem for \(r,s\ge0\).

The finite-tensor approximation step is legitimate because
\(\widehat{\otimes}_{\varepsilon,s,N}X\) is defined as the completion of
\(\otimes^{s,N}X\). The estimate for a finite tensor supported in \(E\)
uses the equality of the injective symmetric norm with the supremum over
\(\Ball_{E^*}\), which follows from Hahn--Banach extension from \(E\) to \(X\).

\section*{Novelty Check and Limitations}

Local duplicate checks in the registry, solution, attempt, and proof-gap
indexes, together with the target-pool scan, found no previous packet for
arXiv:1602.07090 or for Question 4.3.
External searches on 2026-06-19 for exact phrases including
``injective symmetric tensor product'', ``almost square'', ``ASQ'',
``\(\widehat{\otimes}_{\varepsilon,s,N}\)'', and ``1602.07090'' surfaced
the source paper but no later explicit answer. This packet should therefore
be treated as a plausible new affirmative solution, subject to human
literature review.

\section*{Human Review Recommendation}

Recommended review focus: verify Lemma \ref{lem:dual}, the coefficient-safe
finite tensor estimate, and the completion argument. If these pass, the
packet gives a full affirmative answer to Question 4.3 of the source paper.

\begin{thebibliography}{9}

\bibitem{LLRZ}
J. Langemets, V. Lima, and A. Rueda Zoca,
\emph{Almost square and octahedral norms in tensor products of Banach spaces},
arXiv:1602.07090, 2016.

\bibitem{ALL}
T. A. Abrahamsen, J. Langemets, and V. Lima,
\emph{Almost square Banach spaces},
J. Math. Anal. Appl. 434 (2016), 1549--1565.

\end{thebibliography}

\end{document}
